\documentclass{beamer}
\usepackage{amsmath}
\usepackage{gvv}

\title{Question 12.14}
\author{AI25BTECH11040 - Vivaan Parashar}
\date{\today}

\begin{document}

\frame{\titlepage}

\begin{frame}
    \frametitle{Question: }
    One of the eigenvalues of a 3$\times$3 matrix $\vec{M}$ is 3. If the determinant of the matrix $\vec{M}$ is 24 and the trace is 9, then the smallest eigenvalue of the matrix $\vec{M}^{-1}$ is
    \begin{multicols}{2}
        \begin{enumerate}[label=\alph*)]
            \item 1/8
            \item 1/4
            \item 1/3
            \item 1/2
        \end{enumerate}
    \end{multicols}
\end{frame}

\begin{frame}
    \frametitle{Solution: }
    Let the eigenvalues of the matrix $\vec{M}$ be $\lambda_1$, $\lambda_2$, and $\lambda_3$. Given that one of the eigenvalues is 3, we can denote:
    \[\lambda_1 = 3\]
    The determinant of the matrix $\vec{M}$ is given by the product of its eigenvalues:
    \[\text{det}(\vec{M}) = \lambda_1 \cdot \lambda_2 \cdot \lambda_3 = 24\]
    This is because the characteristic polynomial (for a 3$\times$3 matrix) is:
    \[P(\lambda) = \det\left|\vec{M}-\lambda \vec{I}_3\right| = -(\lambda - \lambda_1)(\lambda - \lambda_2)(\lambda - \lambda_3)\]
    Clearly setting $\lambda = 0$ gives the determinant as the product of eigenvalues.
    Substituting the known eigenvalue:
    \[3 \cdot \lambda_2 \cdot \lambda_3 = 24\]
    \begin{align}
        \lambda_2 \cdot \lambda_3 = \frac{24}{3} = 8
    \end{align}
\end{frame}
\begin{frame}
    The trace of the matrix $\vec{M}$ is given by the sum of its eigenvalues (as can be derived from the characteristic polynomial):
    \[\text{tr}(\vec{M}) = \lambda_1 + \lambda_2 + \lambda_3 = 9\]
    Substituting the known eigenvalue:
    \[3 + \lambda_2 + \lambda_3 = 9\]
    \[\lambda_2 + \lambda_3 = 6\]

    Now we have a system of equations:
    \begin{align}
        \lambda_2 + \lambda_3     & = 6 \\
        \lambda_2 \cdot \lambda_3 & = 8
    \end{align}
    We can solve this system by treating $\lambda_2$ and $\lambda_3$ as the roots of the quadratic equation:
    \[x^2 - (\lambda_2 + \lambda_3)x + \lambda_2 \cdot \lambda_3 = 0\]
    Substituting the values from equations (3) and (4):
    \[x^2 - 6x + 8 = 0\]
\end{frame}
\begin{frame}
    This gives us two solutions:
    \begin{align}
        x_1 & = 4 \\
        x_2 & = 2
    \end{align}
    Thus, the eigenvalues of the matrix $\vec{M}$ are:
    \[\lambda_1 = 3, \quad \lambda_2 = 4, \quad \lambda_3 = 2\]
    The eigenvalues of the inverse matrix $\vec{M}^{-1}$ are the reciprocals of the eigenvalues of $\vec{M}$:
    \[\mu_1 = \frac{1}{\lambda_1} = \frac{1}{3}, \quad \mu_2 = \frac{1}{\lambda_2} = \frac{1}{4}, \quad \mu_3 = \frac{1}{\lambda_3} = \frac{1}{2}\]
    The smallest eigenvalue of the matrix $\vec{M}^{-1}$ is:
    \[\min\left(\frac{1}{3}, \frac{1}{4}, \frac{1}{2}\right) = \frac{1}{4}\]
    Thus, the smallest eigenvalue of the matrix $\vec{M}^{-1}$ is:
    \[\boxed{\frac{1}{4}}\]
\end{frame}

\end{document}
